\section*{Bab \ref{ch:g11:seismic-structure} --- Gelombang Seismik Menyingkap Bagian Dalam Bumi}

\begin{solution}{exo:g11:seismic-structure:1}
\omterm{def:g11:seismic-structure:body-waves}{Gelombang P} memampatkan dan meregangkan batuan, yang bergerak searah dengan arah
rambatnya; \omterm{def:g11:seismic-structure:body-waves}{gelombang S} menggesernya, dan batuan bergerak tegak lurus pada arah
itu. \omterm{def:g11:seismic-structure:body-waves}{Gelombang P} menyeberangi zat padat dan zat cair; \omterm{def:g11:seismic-structure:body-waves}{gelombang S} hanya
menyeberangi zat padat. \omterm{def:g11:seismic-structure:body-waves}{Gelombang P} direkam lebih dahulu, karena gelombang itu
yang lebih cepat di antara keduanya.
\end{solution}

\begin{solution}{exo:g11:seismic-structure:2}
$40 \times 111.2 = \qty{4448}{km}$, sekitar \qty{4450}{km}. $8900/111.2
\approx \ang{80}$.
\end{solution}

\begin{solution}{exo:g11:seismic-structure:3}
\qty{456}{s} sama dengan 7~menit 36~detik; \qty{822}{s} sama dengan 13~menit
42~detik. Selang S--P adalah $822 - 456 = \qty{366}{s}$, yaitu 6~menit
6~detik.
\end{solution}

\begin{solution}{exo:g11:seismic-structure:4}
$\sin i_c = 6.0/8.0 = 0.75$, jadi $i_c \approx \ang{48.6}$.
\end{solution}

\begin{solution}{exo:g11:seismic-structure:5}
Tebal inti luar adalah $5150 - 2891 = \qty{2259}{km}$. Jari-jari inti adalah
$6371 - 2891 = \qty{3480}{km}$, dan $3480/6371 \approx 0.55$: sedikit lebih
dari setengah jari-jari Bumi.
\end{solution}

\begin{solution}{exo:g11:seismic-structure:6}
6~menit 6~detik sama dengan \qty{366}{s}: gempa bumi itu berjarak \ang{40}, yaitu
$40 \times 111.2 \approx \qty{4450}{km}$. Satu selang saja hanya memberikan
jarak: episentrum dapat berada di mana saja pada lingkaran berjari-jari itu di
sekeliling stasiun. Diperlukan dua stasiun lagi untuk menemukan satu-satunya
titik yang berjarak tepat dari ketiganya.
\end{solution}

\begin{solution}{exo:g11:seismic-structure:7}
$i_c \approx \ang{48.6}$ dan $\cos i_c \approx 0.661$. Gelombang langsung: $200/6.0
\approx \qty{33.3}{s}$. \omterm{ex:g11:seismic-structure:sp}{Gelombang kepala}: $200/8.0 + 2 \times 35 \times 0.661/6.0
\approx 25.0 + 7.7 = \qty{32.7}{s}$. \omterm{ex:g11:seismic-structure:sp}{Gelombang kepala} tiba lebih dahulu, sekitar
\qty{0.6}{s} lebih awal.
\end{solution}

\begin{solution}{exo:g11:seismic-structure:8}
$\tan i_c \approx 1.134$: jarak kritisnya $2 \times 35 \times 1.134
\approx \qty{79}{km}$. \omterm{prop:g11:seismic-structure:crossover}{Jarak silangnya} $2 \times 35 \times
\sqrt{14.0/2.0} = 70\sqrt 7 \approx \qty{185}{km}$. Stasiun pada
\cref{exo:g11:seismic-structure:7} terletak pada \qty{200}{km}, di luar \omterm{prop:g11:seismic-structure:crossover}{jarak
silang}: \omterm{ex:g11:seismic-structure:sp}{gelombang kepala} tiba lebih dahulu, seperti yang ditemukan.
\end{solution}

\begin{solution}{exo:g11:seismic-structure:9}
$\sin i_c = 6.5/8.0 = 0.8125$, $i_c \approx \ang{54.3}$, dan $\cos i_c \approx
0.583$. \omterm{prop:g11:seismic-structure:head-wave}{Waktu intersep}: $2 \times 6 \times 0.583/6.5 \approx \qty{1.1}{s}$.
\omterm{prop:g11:seismic-structure:crossover}{Jarak silang}: $12\sqrt{14.5/1.5} \approx \qty{37}{km}$. Di bawah benua \omterm{prop:g11:seismic-structure:crossover}{jarak
silang} terletak di dekat \qty{165}{km}, sehingga stasiun yang berjarak beberapa
puluh kilometer satu sama lain menangkap kedua garis itu; di bawah samudra
\omterm{ex:g11:seismic-structure:sp}{gelombang kepala} sudah tiba lebih dahulu di luar sekitar \qty{37}{km}, dan
garis gelombang langsung hanya sepanjang beberapa puluh kilometer: stasiun harus
berjarak beberapa kilometer saja satu sama lain untuk menggambarkannya.
\end{solution}

\begin{solution}{exo:g11:seismic-structure:10}
$\sin r = (8.06/13.7)\sin\ang{30} \approx 0.294$, jadi $r \approx \ang{17.1}$:
sinar itu dibelokkan mendekati garis tegak lurus, yaitu ke arah pusat Bumi.
Karena itu sinar tersebut menyeberangi inti melalui lintasan panjang yang dekat
dengan bagian tengahnya, dan keluar kembali jauh di sisi seberang: sinar P yang
masuk ke inti muncul jauh lebih jauh dari gempa bumi daripada sinar yang nyaris
tidak mengenainya.
\end{solution}

\begin{solution}{exo:g11:seismic-structure:11}
$12\,742/13.7 \approx \qty{930}{s}$, sekitar 15.5~menit, padahal gelombang itu
memerlukan \qty{1210}{s}, sekitar 20~menit: terlambat sekitar 280~detik. Bahkan
dengan kecepatan tertinggi di mantel gelombang itu akan tiba jauh lebih awal,
jadi bagian tengah Bumi pasti memperlambatnya. Oldham menyimpulkan bahwa Bumi
mempunyai inti pusat yang di dalamnya \omterm{def:g11:seismic-structure:body-waves}{gelombang P} merambat lebih lambat daripada
di mantel.
\end{solution}

\begin{solution}{exo:g11:seismic-structure:12}
(1) Menurut proposisi itu, $t_i = 2h\cos i_c/v_1$, dan $\cos i_c = \sqrt{1 -
v_1^2/v_2^2} = \sqrt{v_2^2 - v_1^2}/v_2$. Jadi
\[
  t_i = \frac{2h\sqrt{v_2^2 - v_1^2}}{v_1 v_2},
  \qquad\text{sehingga}\qquad
  h = \frac{t_i\,v_1 v_2}{2\sqrt{v_2^2 - v_1^2}}.
\]
(2) $h = 6.23 \times 6.2 \times 8.1/\bigl(2\sqrt{65.61 -
38.44}\bigr) \approx 312.9/10.43 \approx \qty{30.0}{km}$.
\end{solution}

\begin{solution}{exo:g11:seismic-structure:13}
(1) $(3480/6371)^3 \approx 0.163$: inti merupakan sekitar \qty{16}{\percent}
volume. (2) Massa jenis rata-rata adalah rata-rata yang diberi bobot volume: $5.52 = 0.163
\rho_c + 0.837 \times 4.45$, jadi $\rho_c \approx (5.52 - 3.72)/0.163 \approx
\qty{11.0}{g/cm^3}$. (3) Nilai ini lebih besar daripada \qty{9.9}{g/cm^3},
sebagaimana mestinya, karena massa jenis inti bertambah ke arah pusat (sampai
\qty{13.1}{g/cm^3}); nilai itu sekitar empat kali massa jenis kerak,
\qty{2.6}{g/cm^3}. Inti terbuat dari bahan yang jauh lebih rapat daripada
batuan di permukaan.
\end{solution}

\begin{solution}{exo:g11:seismic-structure:14}
(1) $105 \times 111.2 \approx \qty{11700}{km}$ dan $143 \times 111.2 \approx
\qty{15900}{km}$. (2) $(\cos\ang{105} - \cos\ang{143})/2 \approx (-0.259 +
0.799)/2 \approx 0.27$: sekitar \qty{27}{\percent} permukaan, sebuah cincin di
sekeliling titik yang berseberangan dengan gempa bumi. (3) \omterm{def:g11:seismic-structure:shadow-zone}{Zona bayangan} adalah
rentang jarak dari episentrum: zona itu berpindah bersama gempa buminya. Setiap
gempa bumi mempunyai cincinnya sendiri.
\end{solution}

\begin{solution}{exo:g11:seismic-structure:15}
(1) $h = 9.5 \times 6.3 \times 8.1/\bigl(2\sqrt{65.61 - 39.69}\bigr) \approx
484.8/10.18 \approx \qty{48}{km}$. (2) $\sin i_c = 6.3/8.1 \approx 0.778$,
$\cos i_c \approx 0.629$, dan $t_i = 2 \times 30 \times 0.629/6.3 \approx
\qty{6.0}{s}$. (3) Kerak di bawah jajaran pegunungan itu setebal sekitar
\qty{48}{km}, lebih tebal daripada rata-rata benua yang sekitar \qty{41}{km}
dan dekat dengan rata-rata sabuk pegunungan: pegunungan yang tinggi berdiri di
atas kerak yang tebal.
\end{solution}

\begin{solution}{pb:g11:seismic-structure:1}
\textbf{1.} Titik-titik Pg terletak pada satu garis lurus; titik-titik Pn
terletak pada garis lurus lain yang kurang curam.

\textbf{2.} Gempa bumi itu berada pada kedalaman yang sangat dangkal di bawah
episentrum: stasiun pada jarak nol akan langsung merekam gelombangnya.

\textbf{3.} $(32.3 - 3.3)/(200 - 20) = 29.0/180 \approx \num{0.161}$~\unit{s}
per kilometer.

\textbf{4.} $v_1 = 180/29.0 \approx \qty{6.2}{km/s}$.

\textbf{5.} $3.3 + 280 \times 0.161 \approx \qty{48.4}{s}$.

\textbf{6.} $(43.2 - 18.6)/(300 - 100) = 24.6/200 = \num{0.123}$~\unit{s}
per kilometer, jadi $v_2 = 200/24.6 \approx \qty{8.13}{km/s}$.

\textbf{7.} Nilai itu dekat dengan \qty{8.0}{km/s}: di bawah kerak terletak
batuan mantel paling atas.

\textbf{8.} $\sin i_c = 6.21/8.13 \approx 0.763$, jadi $i_c \approx
\ang{49.8}$.

\textbf{9.} \omterm{ex:g11:seismic-structure:sp}{Gelombang kepala} turun dengan sudut $i_c$ dan naik kembali dengan
sudut yang sama. Titik pertama tempat gelombang itu dapat kembali ke permukaan
terletak sejauh $h\tan i_c$ dari titik tempat ia mencapai \omterm{def:g11:seismic-structure:moho}{Moho}, yang sendiri
berjarak $h\tan i_c$ dari sumber: tidak ada stasiun yang lebih dekat daripada
$2h\tan i_c$ yang dapat menerimanya.

\textbf{10.} Seismometer merekam setiap gelombang yang mencapainya, tidak hanya
yang pertama; gelombang Pn tiba sesudah gelombang Pg tetapi tetap tampak
sebagai gelombang tersendiri pada rekaman.

\textbf{11.} Pada \qty{160}{km}, Pg (\qty{25.8}{s}) tiba sebelum Pn
(\qty{26.0}{s}); pada \qty{180}{km}, Pn (\qty{28.5}{s}) tiba sebelum Pg
(\qty{29.0}{s}). Perubahan itu terjadi di antara kedua stasiun ini.

\textbf{12.} Pg: $t = 3.3 + 0.1611\,(x - 20) \approx 0.08 + 0.1611\,x$. Pn:
$t = 18.6 + 0.123\,(x - 100) = 6.30 + 0.123\,x$ ($x$ dalam kilometer, $t$ dalam
detik).

\textbf{13.} $0.08 + 0.1611\,x = 6.30 + 0.123\,x$ memberikan $x \approx
6.22/0.0381 \approx \qty{163}{km}$.

\textbf{14.} Pn, sesudah $6.30 + 0.123 \times 250 \approx \qty{37.1}{s}$.

\textbf{15.} $t_i \approx \qty{6.3}{s}$.

\textbf{16.} $\cos i_c \approx 0.646$, jadi $h = 6.21 \times 6.3/(2 \times
0.646) \approx \qty{30.3}{km}$.

\textbf{17.} $x_c = 2 \times 30.3 \times \sqrt{14.34/1.92} \approx 60.6 \times
2.73 \approx \qty{165}{km}$, selisihnya kurang dari \qty{2}{km} dengan
\qty{163}{km} pada pertanyaan~13.

\textbf{18.} $\tan i_c \approx 1.18$: $2 \times 30.3 \times 1.18 \approx
\qty{72}{km}$. Stasiun terdekat yang merekam Pn terletak pada \qty{80}{km}, dan
stasiun pada \qty{60}{km} tidak merekamnya: keduanya sesuai.

\textbf{19.} Waktu dibaca sampai \qty{0.1}{s}, dan kesalahan sebesar
\qty{0.1}{s} pada $t_i$ saja mengubah $h$ sekitar \qty{0.5}{km}; kecepatan,
yang diambil dari kemiringan, juga tidak pasti, dan gempa bumi itu tidak tepat
berada di permukaan. Kerak juga tidak seragam sempurna. Ketelitian satu atau dua
kilometer adalah yang paling tinggi yang dimungkinkan oleh data.

\textbf{20.} \textbf{\omterm{def:g11:seismic-structure:moho}{Moho} terletak sekitar \qty{30}{km} di bawah wilayah ini}:
kerak di sini tergolong tipis untuk sebuah benua, yang keraknya rata-rata
setebal sekitar \qty{41}{km}; \qty{30}{km} khas bagi wilayah tempat kerak telah
teregang.
\end{solution}
